\section{Experimental methodology}
\label{sec:methodology}
The study follows a gated progression:
\begin{equation}
\begin{split}
&\text{plant validation}\rightarrow\text{LPV representation validation}
\rightarrow\text{oracle model comparison}\\
&\rightarrow\text{oracle controller comparison}
\rightarrow\text{nonlinear-plant escalation}
\rightarrow\text{federated identification}.
\end{split}
\end{equation}
The early stages deliberately remove learning imperfections. This allows a
negative result to reject the structural premise before a federated algorithm is
developed. Results are first obtained for one deterministic development seed and
then repeated for ten independently sampled fleets and test-trajectory variants.

\subsection{Plant hierarchy}
The first plant is the speed-dependent linear bicycle model from
\cref{sec:benchmark}. It is LTI at a fixed speed but LPV across a drive:
\begin{equation}
 \dot x=A(v_x,\vartheta)x+B(v_x,\vartheta)\delta.
\end{equation}
It supports a clean test of the scheduling--specialization decomposition. After
this oracle test, the plant is replaced by a nonlinear bicycle model with
saturating hyperbolic-tangent tire forces. The identified models remain linear
LTI or LPV models. Thus, the second plant tests whether the compact abstraction
survives deliberate model mismatch.

\subsection{Compared model architectures}
\begin{table}[ht]
\centering
\caption{Core oracle model architectures.}
\label{tab:methods}
\begin{tabular}{llllr}
\toprule
Method & Scheduling & Specialization & Interpretation & Families\\
\midrule
M1: Global LTI & none & none & ignores both factors & 1\\
M2: Clustered LTI & none & oracle family & structure only & 3\\
M3: Global LPV & continuous & none & scheduling only & 1\\
M4: Clustered LPV & continuous & oracle family & both factors & 3\\
M5: Gridded LTI & discrete & family and speed & complex oracle & 15\\
\bottomrule
\end{tabular}
\end{table}
All methods receive equivalent data support. Controller comparisons use common
state and input weights. Method-specific tuning is prohibited unless it is
explicitly reported as a separate sensitivity analysis.

\subsection{Training and test separation}
Frozen-speed identification uses
$v_x\in\{12,16,20,24,28\}\,\mathrm{m\,s^{-1}}$. Interpolation tests include
$v_x\in\{14,18,22,26\}\,\mathrm{m\,s^{-1}}$. Varying-speed training trajectories
combine smooth speed profiles and lateral maneuvers independently so that speed
and maneuver identity are not confounded. Test data include unseen speeds, unseen
speed sequences, unseen lateral timing and amplitude, and unseen speed--maneuver
combinations.

\subsection{Common metrics}
Model metrics include one-step and free-run errors in sideslip and yaw rate, and
matrix errors
\begin{equation}
 E_A(v)=\lVert A_i(v)-\hat A(v)\rVert_F,\qquad
 E_B(v)=\lVert B_i(v)-\hat B(v)\rVert_F.
\end{equation}
Closed-loop metrics include yaw-rate tracking RMSE, sideslip RMS, steering RMS,
peak steering, peak steering rate, and the frozen-grid closed-loop spectral
radius. Every main result reports fleet mean, dispersion, worst-family result,
and family gap
\begin{equation}
 \Delta J_{\mathrm{family}}=\max_c J_c-\min_c J_c.
\end{equation}

\subsection{Experiment documentation contract}
Every experiment below is documented before execution using five fields:
objective, predeclared expectation, setup and evidence, decision gate, and
results/interpretation. The latter remains marked as pending until the experiment
has been run. Configuration, random seeds, output paths, and the producing Git
commit will be added with each result. Failed experiments and changed decisions
remain in this development report even if they are omitted from the final paper.

